\section{同一粗化目标的标量与向量残差控制}
\label{est:partition-controls}
\subsection{一个标量代理与两个同积分的梯度代理}

使用前文的精确积分加残差抽样构造，令标量代理为
\begin{equation*}
 a^{\mathrm K}_v=\ell_{C(v)},\qquad
 I_P(a^{\mathrm K})=K_{\mathcal C}.
\end{equation*}
在 $A\sim P$ 下，相应无偏数值估计量为
\begin{equation*}
 \widehat K_{\mathcal C}
 =K_{\mathcal C}+\ell_A-\ell_{C(A)}.
\end{equation*}
它的均值由链式分解得到完整 $K$，方差为 $\Var_P(\ell_A-\ell_{C(A)})$。桶内条件形状相同时，残差逐点为零，估计量成为确定性的完整 KL。

梯度积分的目标函数为 $\ell_v\psi_v$。同一个粗粒度梯度允许两种可精确积分的代理：
\begin{align*}
 a^{\mathrm{G,const}}_v&=\ell_{C(v)}\bar\psi_{C(v)},\\
 a^{\mathrm{G,score}}_v&=\ell_{C(v)}\psi_v.
\end{align*}
两者逐点不同，但精确积分相同：
\begin{equation*}
 I_P(a^{\mathrm{G,const}})
 =\sum_CP_C\ell_C\bar \psi_C
 =I_P(a^{\mathrm{G,score}})=g_{\mathcal C}.
\end{equation*}
因此它们产生两个不同的无偏梯度估计量：
\begin{align}
 \widehat g_{\mathcal C}^{\mathrm{const}}
 &=g_{\mathcal C}+\ell_A\psi_A-\ell_{C(A)}\bar\psi_{C(A)},
 \label{inf:eq:kl-partition-constant-estimator}\\
 \widehat g_{\mathcal C}^{\mathrm{score}}
 &=g_{\mathcal C}+(\ell_A-\ell_{C(A)})\psi_A.
 \label{inf:eq:kl-partition-score-estimator}
\end{align}
两种代理之差为
\begin{equation*}
 \chi_v=\ell_{C(v)}(\psi_v-\bar\psi_{C(v)}),
 \qquad \E_{P^C}[\chi_A]=0.
\end{equation*}
所以改变这两个代理之间的线性组合，只改变桶内零均值波动。它不改变粗粒度精确积分，也不改变每个桶内残差的条件均值。在前文控制矩阵与 Gram 二次型框架中，$\chi_A$ 就是一列已知均值为零的控制特征；其最优系数由残差与该特征的协方差决定，取系数一不自动保证降低方差。

\subsection{头部 singleton 与尾部单桶的特例}

令 $S$ 为冻结的 teacher Top-$J$ 支持，定义 $P_S=\sum_{v\in S}P(v)$、$Q_S=\sum_{v\in S}Q(v)$、$\tau_P=1-P_S$、$\tau_Q=1-Q_S$ 和 $J_A=\mathbf1\{A\notin S\}$，本节头尾公式假设 $0<P_S,Q_S<1$。头部积分为
\begin{equation*}
 H=\sum_{v\in S}P(v)\ell_v,\qquad h_g=\sum_{v\in S}P(v)\ell_v\psi_v.
\end{equation*}
令 $\mathcal C=\{\{v\}:v\in S\}\cup\{\mathcal V\setminus S\}$。头部单元满足 $\ell_{\{v\}}=\ell_v$、$\bar\psi_{\{v\}}=\psi_v$；尾部单元满足
\begin{equation*}
 \ell_{\mathrm{mass}}=\log\frac{\tau_P}{\tau_Q},\qquad
 \bar \psi_\tail=\nabla\log\tau_P,\qquad
 b_{\mathrm{TA}}=\ell_{\mathrm{mass}}\bar \psi_\tail.
\end{equation*}
因此 $K_{\mathcal C}=H+\tau_P\ell_{\mathrm{mass}}$ 是 Tail-aware KL，且 $g_{\mathcal C}=h_g+\tau_Pb_{\mathrm{TA}}$。三个残差构造分别成为
\begin{align*}
 \widehat K_{\mathcal C}
 &=H+\tau_P\ell_{\mathrm{mass}}+J_A(\ell_A-\ell_{\mathrm{mass}}),\\
 \widehat g_{\mathcal C}^{\mathrm{const}}
 &=h_g+\tau_Pb_{\mathrm{TA}}+J_A(\ell_A\psi_A-b_{\mathrm{TA}}),\\
 \widehat g_{\mathcal C}^{\mathrm{score}}
 &=h_g+\tau_Pb_{\mathrm{TA}}+J_A(\ell_A-\ell_{\mathrm{mass}})\psi_A.
\end{align*}
第一式是标量残差估计量；第二式是 ResOPD；第三式是 SC-TA。HETS 使用头部精确代理和零尾部代理，因而是另一个代理选择。上述分类从积分目标和代理函数同时确定估计量，不依赖算法名称。

\begin{center}
\begin{tabular}{lll}
\toprule
目标与方法 & 头部近似 $a_v$ & 尾部近似 $a_v$\\
\midrule
数值 HETS & $\ell_v$ & $0$\\
数值桶残差校正 & $\ell_v$ & $\ell_{\mathrm{mass}}$\\
梯度 HETS & $\ell_v\psi_v$ & $0$\\
梯度 ResOPD & $\ell_v\psi_v$ & $b_{\mathrm{TA}}$\\
梯度 SC-TA & $\ell_v\psi_v$ & $\ell_{\mathrm{mass}}\psi_v$\\
\bottomrule
\end{tabular}
\end{center}
表中每一行都使用式~\eqref{est:eq:kl-master-estimator}；确定积分项是该行的 $I_P(a)$，采样项始终是 $f_A-a_A$。数值 Tail-aware 与梯度粗化则只返回各自行的确定积分项，没有补回残差。

\begin{proposition}[标量精确时，常向量代理仍可有方差]
若 $P^{\tail}=Q^{\tail}$，则上述标量残差估计量和 SC-TA 都是确定性的完整目标，而 ResOPD 的协方差为
\begin{equation*}
 \operatorname{Cov}(\widehat g^{\mathrm{ResOPD}})
 =\tau_P\ell_{\mathrm{mass}}^2\operatorname{Cov}_{P^{\tail}}(\psi_A).
\end{equation*}
\end{proposition}

\begin{proof}
桶内条件分布相同意味着尾部 $\ell_v=\ell_{\mathrm{mass}}$。标量残差和 SC-TA 的随机残差于是逐点为零。ResOPD 则仍有随机残差
\begin{equation*}
 J_A\ell_{\mathrm{mass}}(\psi_A-\bar \psi_\tail).
\end{equation*}
它在尾部条件下均值为零，在头部为零，因此总体均值为零，二阶矩恰好是命题中的协方差。
\end{proof}

例如取 $P=(0.4,0.2,0.4)$、$Q=(0.7,0.1,0.2)$ 和 $S=\{1\}$，则两个尾部条件分布均为 $(1/3,2/3)$，且 $\tau_P=0.6$、$\ell_{\mathrm{mass}}=\log2$。在完整 softmax logits 参数化下，$\psi_v=e_v-P$，尾部 score 的协方差迹为 $1-(1/3)^2-(2/3)^2=4/9$。ResOPD 的方差因此为 $\frac4{15}(\log2)^2>0$；标量残差估计量与 SC-TA 的方差均为零。这个例子直接区分了可精确积分的 KL 数值、粗粒度真实梯度以及不同梯度代理留下的抽样波动。
